\section{Introduction}

Satellite images contain useful information about land use and land cover. This
makes automatic classification of such images a natural application for deep
learning. In this project, the EuroSAT RGB dataset is used
\cite{helber2019eurosat,eurosatgithub}.

The project compares two model approaches on the same dataset. Model 1 is a CNN
trained from scratch. Its purpose is to show the full training process, from
architecture design to optimizer choice and evaluation. Model 2 is a transfer
learning approach, in which an ImageNet-pretrained convolutional backbone is
fine-tuned on EuroSAT. To avoid drawing conclusions from a single architecture,
the transfer learning approach is evaluated with two structurally different
backbones: the parameter-efficient EfficientNet-B0
\cite{tan2019efficientnet,pytorch} and the larger deep-residual ResNet-50
\cite{he2016resnet}. Comparing two backbones shows whether any advantage over
the scratch CNN comes from transfer learning in general rather than from one
specific model, and lets us study how pretrained features affect accuracy and
data efficiency across architectures.

Two tasks are considered on the same dataset:

\begin{itemize}
  \item a 10-class classification task using the original EuroSAT classes,
  \item a binary classification task: vegetation vs. non-vegetation.
\end{itemize}

The same train/validation/test split was used for both tasks. This is important
because otherwise differences in the results could also come from different
data partitions. The hyperparameter experiments used only the training and
validation splits. The test split was evaluated only for the final selected
configurations.
